% -----------------------------------------------------------
\section{Seal}
% -----------------------------------------------------------
	\label{SEAL}
	
% % ----------------------
% \subsection{See also}
% % ----------------------

% % ----------------------
% \subsection{1828 American Dictionary of the English Language} % *1828
% % ----------------------
	% \label{SEAL-1828}

	% \begin{compactitem}
		% \item
	% \end{compactitem}

% ----------------------
\subsection{Oxford English Dictionary} % *OED
% ----------------------
	\label{SEAL-OED}

	\begin{compactitem}
		\item[I] To attest by a seal.
		\item[I.1.a] \textit{transitive}. To place a seal upon (a document) as evidence of genuineness, or as a mark of authoritative ratification or approval.
		\item[I.1.b] \textit{figurative}. To authenticate or attest solemnly by some act compared to the affixing of a seal.
		\item[I.1.c] To conclude, ratify, render binding (an agreement, etc.) by affixing the seals of the parties to the instrument. Also \textit{figurative}, to ratify or clinch (a bargain) by some ceremonial act.
		\item[I.1.d] To grant (a charter, etc.) under one's seal; \textit{figurative} to give (testimony, a promise, etc.) in an authoritative manner or with solemn pledges of good faith.
		\item[I.1.e] To impose (an obligation, a penalty) on a person in a binding manner.
		\item[I.1.f] \textit{figurative}. Of a thing or act: To attest or ratify as a seal does; to be a `seal' of.
		\item[I.1.g] To decide irrevocably (the fate of a person or thing); to complete and place beyond dispute or reversal (a victory, defeat, etc.).
		\item[I.2.a] To mark by a seal as reserved for a particular destination. Chiefly \textit{figurative}, esp. in certain uses of New Testament origin: To designate, set apart, assign to another person or bind \textit{together}, by an inviolable token or pledge.
		\item[I.2.b] In allusions to Revelation vii. 5--8.
		\item[I.2.c] Among the Mormons, to set apart (a woman) by a solemn ceremony to a man as one of his `spiritual wives'.
		\item[I.3] To impress a seal upon (weights or measures) to indicate that their correctness has been tested by municipal or other lawful authority. Also, to place an official stamp on (merchandise, e.g. pieces of cloth) to certify that it is of standard measure or quality.
		\item[I.4.a] \textit{intransitive}. (Cf. 1 \textit{absol}.) To set one's seal (to a document). Also \textit{spec}. to set one's seal to or execute a promissory note; to become security \textit{for} a person. Also to \textbf{\textit{seal under}}. \textit{Obsolete}.
		\item[I.4.b] \textit{figurative}. To assent, lend one's support or authority to a statement or proposal. \textit{Obsolete}.
		\item[I.4.c] \textit{figurative} ? To make peace. \textit{Obsolete}.
		\item[II] To fasten with or as with a seal.
		\item[II.5.a] \textit{transitive}. To fasten (a folded letter or other document) with melted wax or some other plastic material and impress a seal upon this, so that opening is impossible unless the seal is broken.
			\begin{compactitem}
				\item Dating back to around 1225.
				\item Note that, as it says in TDNT, the earliest seals were clay, not wax.
			\end{compactitem}
		\item[II.5.b] Said of the signet itself.
		\item[II.5.c] To stamp the wax fastening (a letter) \textit{with} something substituted for a seal.
		\item[II.5.d] To fasten up (a letter, a parcel) with sealing-wax, a wafer, gum, or the like.
		\item[II.6.a] To place a seal upon the opening of (a door, a chest, etc.) for security. Also with \textit{up}.
		\item[II.6.b] In figurative phrases, esp. \textbf{\textit{to seal (}a person's\textit{) lips}}, to bind or constrain to silence or secrecy; \textbf{\textit{to seal (}a person's\textit{) eyes} or \textit{ears}}, to render blind or deaf, also to restrain from looking or listening. Also with \textit{up}.
		\item[II.7.a] To place in a receptacle secured by a seal.
		\item[II.7.b] Chess. \textbf{\textit{to seal a move}}. Of a player in a match or tournament: To place in a sealed envelope a statement of the move he intends to make when the game is resumed after an adjournment.
		\item[II.8.a] To close (a vessel, an aperture, etc.) securely by placing a coating of wax, cement, or lead, over the orifice, or, in wider sense, by any kind of fastening that must be broken before access can be obtained. \textbf{\textit{to seal hermetically}}: see the adverb.
		\item[II.8.b] \textit{Surgery}. To close up (a wound) with a covering that is not to be removed until healing has taken place.
		\item[II.8.c] To render (a surface of wood, etc.) impervious by the application of a special coating.
		\item[II.8.d] To prevent access to and egress from (an area or space); to close (entrances) for this purpose. Usually with \textit{off}.
		\item[II.9.a] To fasten \textit{on} or \textit{down} with wax or cement.
		\item[II.9.b] \textit{figurative}. To fasten, fix immoveably.
		\item[II.9.c] \textit{intransitive}. To apply wax, etc. in fastening.
		\item[II.10] \textit{Building}. \textit{transitive}. To secure in position in a wall by means of mortar, cement, etc. [So French \textit{sceller}.]
		\item[II.11] \textit{transferred}. To enclose, shut \textit{up} within impenetrable barriers. Also, to confine so as to prevent access or egress, and with \textit{off}.
		\item[III] To mark with or as with a seal.
		\item[III.12.a] \textit{transitive}. To impress (a mark) \textit{upon}.
		\item[III.12.b] \textit{intransitive}. To make marks like those on a seal. \textit{Obsolete}.
	\end{compactitem}
	
% % ----------------------
% \subsection{Restoration Lexicon} % *RL *RESTORATION LEXICON
% % ----------------------
	% \label{SEAL-OED}

	% \begin{compactitem}
		% \item[]
	% \end{compactitem}

% ----------------------
\subsection{Scriptural Usage} % *SCRIPTURES
% ----------------------
	\label{SEAL-SCRIPTURES}

	% % -----
	% \subsubsection{Old Testament} % *OT
	% % -----
		% \label{SEAL-OT}
	
		% % --
		% \paragraph{KJV}
		% % --
			% \label{SEAL-OT-KJV}
	
			% \begin{tabular}{ll}
				% Total occurrences in the OT & 
			% \end{tabular}
			
			% \begin{compactdesc}
				% \item[]
			% \end{compactdesc}
			
		% % --
		% \paragraph{Hebrew Bible}
		% % --
			% \label{SEAL-HEBREW}
		
			% %
			% \subparagraph{\texorpdfstring{\he{sample word}}{}} % paste actual hebrew text here
			% %
				% \label{PAGA} % whatever is in step bible without periods
			
				% \begin{compactdesc}
					% \item[HALOT , PDF ] \textbf{} % Use the 1994 PDF
						% \begin{compactitem}
							% \item[1]
						% \end{compactitem}
					% \item[BDB , PDF ] \textbf{}
						% \begin{compactitem}
							% \item[1]
						% \end{compactitem}
					% \item[DCH PDF ] \textbf{} % don't include the actual page number, they repeat from volume to volume
						% \begin{compactitem}
							% \item[1]
						% \end{compactitem}
				% \end{compactdesc}
				
				% \begin{compactdesc}
					% \item[Some OT uses] \textbf{}
						% \begin{compactdesc}
							% \item[]
						% \end{compactdesc}
				% \end{compactdesc}
			
		% % --
		% \paragraph{Septuagint}
		% % --
			% \label{SEAL-LXX}
		
			% %
			% \subparagraph{\texorpdfstring{\gr{sample word}}{}}
			% %
		
	% -----
	\subsubsection{New Testament} % *NT
	% -----
		\label{SEAL-NT}
	
		% --
		\paragraph{KJV}
		% --
			\label{SEAL-NT-KJV}		
	
			\begin{tabular}{ll}
				Total occurrences in the NT & 42
			\end{tabular}
			
			\begin{compactdesc}
				\item[{\hyperref[EPHESIANS-1-13]{Ephesians 1:13}}]
			\end{compactdesc}
			
		% --
		\paragraph{Greek New Testament} % *GREEK
		% --
			\label{SEAL-GREEK}
		
			%
			\subparagraph{\texorpdfstring{\gr{σφραγίζω}}{}}
			%
				\label{SPHRAGIZO}
			
				\begin{compactdesc}
					\item[BDAG 871a] \textbf{}
						\begin{compactitem}
							\item[1] \textbf{to provide with a seal as a security measure, \textit{seal}} with accusative of the object that is to be secured or fastened by the seal
							\item[2] \textbf{to close something up tight, \textit{seal}}
							\item[3] \textbf{to mark with a seal as a means of identification,\textit{mark, seal}}...so that the mark denoting ownership also carries with it the protection of the owner...\gr{σφραγίσωμεν τοὺς δούλους τοῦ θεοῦ ἐπὶ τῶν μετώπων αὐτῶν} (\hyperref[REVELATIONNT-7-3]{Revelation 7:3})..this forms a basis for understanding the imagery which speaks of those who enter the Christian community as being sealed with or by the Holy Spirit (\hyperref[EPHESIANS-1-13]{Ephesians 1:13}; compare \hyperref[EPHESIANS-4-30]{Ephesians 4:30}). Similarly \gr{θεός, ὁ σφραγισάμενος ἡμᾶς καὶ δοὺς τὸν ἀρραβῶνα τοῦ πνεύματος ἐν ταῖς καρδίαις ἡμῶν} (\hyperref[2CORINTHIANS-1-22]{2 Corinthians 1:22}), but here the context contributes an additional component in the sense of `endue with power from heaven', as plainly in \hyperref[JOHN-6-27]{John 6:27}...
							\item[4] \textbf{to certify that something is so, \textit{attest, certify, acknowledge}} figurative (as a seal does on a document)
							\item[5] \textbf{to seal something for delivery, \textit{seal}}
								\begin{compactitem}
									\item There is a discussion of \hyperref[ROMANS-15-28]{Romans 15:28} here.
								\end{compactitem}
						\end{compactitem}
					\item[CGL 1349b, PDF 1389] \textbf{}
						\begin{compactitem}
							\item[1] mark with a seal (for identification or security), \textbf{seal}---\textit{a document}, \textit{a boulder} (blocking an entrance) (NT); (middle)---\textit{a document}, \textit{the doors of a storeroom}
							\item[2] (figurative) mark with a seal (for authentication or approval); (of God) \textbf{mark with His seal} ---\textit{Christ} (NT); (of a person) \textbf{certify} ---with complementary clause \textit{that something is the case} (NT)
							\item[3] passive (figurative, of persons who have been in a fight) \textbf{be stamped} ---with dative \textit{with marks of injury}
						\end{compactitem}
					\item[LSJ 1742a, PDF 1813] \textbf{}
						\begin{compactitem}
							\item[I.1] \textit{close} or \textit{enclose with a seal}
							\item[I.2] \textit{authenticate} a document \textit{with a seal}
							\item[I.3] \textit{certify} an object \textit{after examination by attaching a seal}
							\item[I.4] Middle, \textit{seal} an article to show that it is pledged
							\item[II] metaphorical senses
							\item[II.1] \textit{close up as if with a seal}
							\item[II.2] \textit{accredit} as an envoy
							\item[II.3] \textit{set a seal of approval upon, confirm}
							\item[II.4] generally, \textit{mark}
							\item[II.5] \textit{set an end} or \textit{limit to}
						\end{compactitem}
				\end{compactdesc}
				
				\begin{tabular}{ll}
					Total occurrences in the NT & 15
				\end{tabular}
				
				\begin{compactdesc}
					\item[Some NT uses] \textbf{}
						\begin{compactdesc}
							\item[{\hyperref[MATTHEW-27-66]{Matthew 27:66}}]
							\item[{\hyperref[JOHN-3-33]{John 3:33}}]
							\item[{\hyperref[JOHN-6:27]{John 6:27}}]
							\item[{\hyperref[ROMANS-15-28]{Romans 15:28}}]
							\item[{\hyperref[2CORINTHIANS-1-22]{2 Corinthians 1:22}}]
							\item[{\hyperref[EPHESIANS-1-13]{Ephesians 1:13}}]
							\item[{\hyperref[EPHESIANS-4-30]{Ephesians 4:30}}]
							\item[{\hyperref[REVELATIONNT-7-3]{Revelation 7:3}}]
							\item[{\hyperref[REVELATIONNT-7-4]{Revelation 7:4}}] \textbf{(2)}
							\item[{\hyperref[REVELATIONNT-7-5]{Revelation 7:5}}]
							\item[{\hyperref[REVELATIONNT-7-8]{Revelation 7:8}}]
							\item[{\hyperref[REVELATIONNT-10-4]{Revelation 10:4}}]
							\item[{\hyperref[REVELATIONNT-20-3]{Revelation 20:3}}]
							\item[{\hyperref[REVELATIONNT-22-10]{Revelation 22:10}}]
						\end{compactdesc}
				\end{compactdesc}
				
			%
			\subparagraph{TDNT: \texorpdfstring{\gr{σφραγίζω}}{}  (7:939--953)}
			%
				\label{SPHRAGIZO-TDNT}
		
	% % -----
	% \subsubsection{Book of Mormon} % *BOM
	% % -----
		% \label{SEAL-BOM}
	
		% \begin{tabular}{ll}
			% Total occurrences in the BOM & 
		% \end{tabular}
		
		% \begin{compactdesc}
			% \item[]
		% \end{compactdesc}
		
	% -----
	\subsubsection{Doctrine and Covenants} % *DC
	% -----
		\label{SEAL-DC}
	
		% \begin{tabular}{ll}
			% Total occurrences in the DC & 
		% \end{tabular}
		
		\begin{compactdesc}
			\item[{\hyperref[DC1-8]{DC 1:8--9}}]
		\end{compactdesc}
		
	% % -----
	% \subsubsection{Pearl of Great Price} % *PGP
	% % -----
		% \label{SEAL-PGP}
	
		% \begin{tabular}{ll}
			% Total occurrences in the PGP & 
		% \end{tabular}
		
		% \begin{compactdesc}
			% \item[]
		% \end{compactdesc}
		
% % ----------------------
% \subsection{Key Phrases} % *KEYPHRASES *PHRASES
% % ----------------------
	% \label{SEAL-PHRASES}

	% \begin{compactdesc}
		% \item[] \textbf{\#times} | list
			% % \begin{compactdesc}
				% % \item[Bible anchor verses] \phantom{.}
					% % \begin{compactdesc}
						% % \item[Romans 5:5] \gr{}
					% % \end{compactdesc}
				% % \item[Commentary] ---
			% % \end{compactdesc}
	% \end{compactdesc}
	
% % ----------------------
% \subsection{Bible Dictionaries} % *DICTIONARIES
% % ----------------------
	% \label{SEAL-BD}

% % ----------------------
% \subsection{Restoration Usage} % *RESTORATION
% % ----------------------
	% \label{SEAL-RESTORATION}

	% % -----
	% \subsubsection{Joseph Smith} % *JS
	% % -----
		% \label{SEAL-JS}
	
		% \begin{compactdesc}
			% \item[{\hyperref[KEY]{Date/Source}}]
		% \end{compactdesc}
		
	% % -----
	% \subsubsection{Temple Ordinances}
	% % -----
		% \label{SEAL-TEMPLE}
	
		% \begin{compactdesc}
			% \item[{\hyperref[KEY]{Date/Source}}]
		% \end{compactdesc}
	
	% % -----
	% \subsubsection{Brigham Young} % *BY
	% % -----
		% \label{SEAL-BY}
	
		% \begin{compactdesc}
			% \item[{\hyperref[KEY]{Date/Source}}]
		% \end{compactdesc}
	
% % ----------------------
% \subsection{Commentary and Studies} % *COMMENTARY *STUDIES
% % ----------------------
	% \label{SEAL-COMMENTARY}

	% % -----
	% \subsubsection{Key Points}
	% % -----
		% \label{SEAL-KEYS}
	
		% \begin{compactitem}
			% \item
		% \end{compactitem}
		
	% % -----
	% \subsubsection{Sample Study}
	% % -----
		% \label{SEAL-study}